\section{Semantic segmentation}
\seclabel{segmentation}

Region classification is a standard technique for semantic segmentation, allowing us to easily apply R-CNN to the PASCAL VOC segmentation challenge.
To facilitate a direct comparison with the current leading semantic segmentation system (called \OtwoP for ``second-order pooling'') \cite{o2p}, we work within their open source framework.
\OtwoP uses CPMC to generate 150 region proposals per image and then predicts the quality of each region, for each class, using support vector regression (SVR).
The high performance of their approach is due to the quality of the CPMC regions and the powerful second-order pooling of multiple feature types (enriched variants of SIFT and LBP).
We also note that Farabet \etal \cite{farabet-pami-13} recently demonstrated good results on several dense scene labeling datasets (not including PASCAL) using a CNN as a multi-scale per-pixel classifier.

We follow \cite{arbelaez2012semantic,o2p} and extend the PASCAL segmentation training set to include the extra annotations made available by Hariharan \etal \cite{hariharan2012inverse}. 
Design decisions and hyperparameters were cross-validated on the VOC 2011 validation set. Final test results were evaluated only once.

\paragraph{CNN features for segmentation.} We evaluate three strategies for computing features on CPMC regions, all of which begin by warping the rectangular window around the region to $227 \times 227$.
The first strategy (\textit{full}) ignores the region's shape and computes CNN features directly on the warped window, exactly as we did for detection.
However, these features ignore the non-rectangular shape of the region.
Two regions might have very similar bounding boxes while having very little overlap.
Therefore, the second strategy (\textit{fg}) computes CNN features only on a region's foreground mask.
We replace the background with the mean input so that background regions are zero after mean subtraction.
The third strategy (\textit{full+fg}) simply concatenates the \textit{full} and \textit{fg} features; our experiments validate their complementarity.
\begin{table}[h!]
\centering
{\small
\begin{tabular}{@{}c|c|c|c|c|c|c}
& \multicolumn{2}{c|}{\textit{full} R-CNN}
& \multicolumn{2}{c|}{\textit{fg} R-CNN}
& \multicolumn{2}{c}{\textit{full+fg} R-CNN}
\\
\hline
\OtwoP \cite{o2p}
& \fc{6} & \fc{7} & \fc{6} & \fc{7} & \fc{6} & \fc{7} \\
\hline
  46.4   % o2p
& 43.0   % full    CNN fc6
& 42.5   % full    CNN fc7
& 43.7    % fg      CNN fc6
& 42.1    % fg      CNN fc7
& \textbf{47.9}   % full+fg CNN fc6
& 45.8    % full+fg CNN fc7
\end{tabular}
}
\caption{\textbf{Segmentation mean accuracy (\%) on VOC 2011 validation.}
Column 1 presents \OtwoP;
2-7 use our CNN pre-trained on ILSVRC 2012.
}
\tablelabel{voc2011val}
\vspace{-1em}
\end{table}
\paragraph{Results on VOC 2011.}
\tableref{voc2011val} shows a summary of our results on the VOC 2011 validation set compared with \OtwoP.
(See \asecref{segperclass} for complete per-category results.)
Within each feature computation strategy, layer \fc{6} always outperforms \fc{7} and the following discussion refers to the \fc{6} features.
The \textit{fg} strategy slightly outperforms \textit{full}, indicating that the masked region shape provides a stronger signal, matching our intuition.
However, \textit{full+fg} achieves an average accuracy of 47.9\%, our best result by a margin of 4.2\% (also modestly outperforming \OtwoP), indicating that the context provided by the \textit{full} features is highly informative even given the \textit{fg} features.
Notably, training the 20 SVRs on our \textit{full+fg} features takes an hour on a single core, compared to 10+ hours for training on \OtwoP features.

In \tableref{voc2011test} we present results on the VOC 2011 test set, comparing our best-performing method, \fc{6} (\textit{full+fg}), against two strong baselines.
Our method achieves the highest segmentation accuracy for 11 out of 21 categories, and the highest overall segmentation accuracy of 47.9\%, averaged across categories (but likely ties with the \OtwoP result under any reasonable margin of error).
Still better performance could likely be achieved by fine-tuning.

\begin{table*}[t!]
\centering
\renewcommand{\arraystretch}{1.2}
\renewcommand{\tabcolsep}{1.2mm}
\resizebox{\linewidth}{!}{
\begin{tabular}{@{}l|c|r*{19}{c}|c@{}}
\textbf{VOC 2011 test} & bg & aero      & bike      & bird      & boat      & bottle     & bus        & car        & cat        & chair      & cow        & table      & dog        & horse      & mbike      & person     & plant      & sheep      & sofa       & train      & tv         & mean \\
\hline
R\&P \cite{arbelaez2012semantic} &
83.4 &
46.8 & 18.9 & 36.6 & 31.2 & 42.7 & 57.3 & 47.4 & 44.1 & \phz8.1 & 39.4 &
\bf{36.1} & 36.3 & 49.5 & 48.3 & 50.7 & 26.3 & 47.2 & 22.1 & 42.0 & 43.2 &
40.8 \\
\OtwoP \cite{o2p} &
\bf{85.4} &
\bf{69.7} & 22.3 & 45.2 & \bf{44.4} & 46.9 & 66.7 & 57.8 & 56.2 & \bf{13.5} & \bf{46.1} &
32.3 & 41.2 & \bf{59.1} & 55.3 & 51.0 & \bf{36.2} & 50.4 & \bf{27.8} & 46.9 & \bf{44.6} &
47.6 \\
\hline
\textbf{ours} (\textit{full+fg} R-CNN \fc{6}) &
84.2 &
66.9 & \bf{23.7} & \bf{58.3} & 37.4 & \bf{55.4} & \bf{73.3} & \bf{58.7} & \bf{56.5} & \phz9.7 & 45.5 &
29.5 & \bf{49.3} & 40.1 & \bf{57.8} & \bf{53.9} & 33.8 & \bf{60.7} & 22.7 & \bf{47.1} & 41.3 &
\bf{47.9} \\
\end{tabular}
}
\caption{\textbf{Segmentation accuracy (\%) on VOC 2011 test.}
%All methods compete in the \textit{comp6} track due to use of outside data from \cite{hariharan2012inverse}.
We compare against two strong baselines: the ``Regions and Parts'' (R\&P) method of~\cite{arbelaez2012semantic} and the second-order pooling (\OtwoP) method of~\cite{o2p}.
Without any fine-tuning, our CNN achieves top segmentation performance, outperforming R\&P and roughly matching \OtwoP.
}
\tablelabel{voc2011test}
\vspace{-1em}
\end{table*}
